% SHARD-ID: RC-04L-ELLIPTIC-CAVITY-CHANNEL
% SHARD-TITLE: The Phantasm VI: Two Arithmetic Cavities with Cusps II. The h-Exit Model, the Renewal Channel, the Rebounds the Arithmetic Does and Does Not Pick, the Arithmetic Metric, the Vacuum, and Funnels
% SHARD-SUMMARY: The one-exit machinery of shards 04i/04j is generalised to h exits on the model side: an ordered Blaschke-Potapov pole removal gives a rational inner matrix function Theta, the model space has dimension deg det Theta = number of resonances with the delay modes included, the defect I - Z^*Z = J^*J has rank at most h, and the renewal channel E(rho) = Z rho Z^* + R(J rho J^*) with a CP rebound map R is CPTP; stationary densities correspond bijectively to stationary exit densities of an h^2 x h^2 renewal instrument, uniqueness is not irreducibility, and the common-modulus theorem holds for fixed resets but fails for exit-dependent rebounds.
% SHARD-SUMMARY: On the one-cusp elliptic cavity the canonical cusp rebound is a pure delay mode (Theta(0) = 0 because c = 1): absorbing, holding law delta_1, never charging the Hasse-Weil sector, so "flux returns where it left" is not the arithmetic rebound; modal Hasse-Weil rebounds are all stationary (Weil) with mean holding time 2 + sqrt 2, exit rate 1 - q^{-1/2}, and Gram spectrum, not weight list; the metrics in which q^{1/4} Z is unitary form a three-parameter cone whose parity-and-reality invariant part is a single ray, the Lax-Phillips metric is outside it (angles 80.1, 65.5, 79.3 degrees), and in the arithmetic metric the Hasse-Weil defect has rank four, so the arithmetic channel needs one exit per zero mode.
% SHARD-SUMMARY: The constant function on a finite-volume regular diagram is a bound state with wave multipliers q^{+-1/2}, never stationary and never leaking: no number of cusps gives the vacuum an exit; a funnel enters the resonance determinant with self-energy C/(qz), turning the constant mode from a bound state at z = q^{-1/2} into an outgoing state at the same value, but a globally constant function cannot be outgoing on a cusp and a funnel at once, and cusp-plus-funnel diagrams are not lattice quotients. Hecke selection of the rebound and the genus/class-number prediction are open.
% SHARD-KEYWORDS: h exits, Blaschke-Potapov, inner matrix function, model space, defect, Gram identity, renewal channel, CP rebound, renewal instrument, fixed reset, common modulus, delay mode, canonical rebound, absorbing, Hasse-Weil sector, arithmetic metric, unitarisability, vacuum, finite volume, funnel, Hecke, H-HECKE, H-CLASS, H-LP
% SHARD-NOTES: none
% SHARD-SCRIPTS: scripts/elliptic_cavity.py

\section{The Phantasm VI: Two Arithmetic Cavities with Cusps II. The $h$-Exit Model, the Renewal Channel, the Arithmetic Metric, and the Vacuum}
\label{sec:elliptic-cavity-channel}

Same lanes and files as \cref{sec:elliptic-cavity-scattering}. Throughout, H-CORE means a finite real
symmetric core with finitely many standard rays and positive couplings (\cref{def:elliptic-cavity}), H-LP is
the energy completion and the repair $D_+ = \mathcal O_+\cap\mathcal O_-^\perp$ of
\cref{thm:cusp-wave-contraction} (inherited, not re-litigated: the model below is proved, its
identification with the geometric wave compression is not), and H-CONTRACTION means a finite-dimensional
$Z$ with $Z^m\to0$ and $I-Z^{*}Z = J^{*}J$.

\subsection{The model with $h$ exits and the renewal channel}

\begin{definition}[The $h$-exit model; rebound map; renewal channel; fixed reset]
\label{def:h-exit-renewal-channel}
For $\mathbf S$ as in \cref{def:elliptic-cavity}, the \emph{inner part} is $\Theta = \eta\,\mathbf S\,B_{\mathrm{bd}}$
with $B_{\mathrm{bd}}$ the ordered right Blaschke--Potapov product removing the bound-state poles and $\eta$ a
constant unitary; the \emph{model space} is $K = H^2(\mathbb C^h)\ominus\Theta H^2(\mathbb C^h)$, the
\emph{compressed shift} $Z = P_KM_w|_K$, and the \emph{exit map} $J:K\to\mathbb C^h$ is defined by
$(I-P_K)(wf) = \Theta\,Jf$. A \emph{rebound map} is a CPTP map $\mathsf R:\mathcal B(\mathbb C^h)\to\mathcal B(K)$;
the \emph{renewal channel} is $\mathcal E_{\mathsf R}(\rdens) = Z\rdens Z^{*}+\mathsf R(J\rdens J^{*})$; its
\emph{holding instrument} is the sequence of completely positive maps $\mathcal M_m = \Psi\mathcal E_0^{m-1}\mathsf R$
on $\mathcal B(\mathbb C^h)$, $\Psi(X) = JXJ^{*}$, $\mathcal E_0(X) = ZXZ^{*}$, with $\mathcal R = \sum_{m\ge1}\mathcal M_m$
the \emph{exit renewal channel}. A \emph{fixed reset} is $\mathsf R(X) = \mathrm{Tr}(X)\,\reb$ for a density $\reb$
on $K$; a \emph{modal} reset has $\reb$ a mixture of normalised eigenvectors of $Z$. The \emph{Hasse--Weil sector}
$K_{\mathrm{HW}}$ of an elliptic cavity is the invariant span of the eigenvectors at the nonzero resonances; the
\emph{delay sector} is the generalised eigenspace at $0$.
\end{definition}

\begin{theorem}[The $h$-exit model: inner part, dimension, defect, Gram identity]
\label{thm:h-exit-model}
\claimstatus{thm:h-exit-model}{proved}
Under H-CORE. Removing a simple matrix pole at $\vartheta$ with residue $R_\vartheta$ by the right factor
$B_{\vartheta,P} = I-P+b_\vartheta P$, $P$ the projection onto $\mathrm{ran}\,R_\vartheta^{*}$, $b_\vartheta(z) = (z-\vartheta)/(1-\bar\vartheta z)$,
and repeating with recomputed residues, gives an ordered minimal product $B_{\mathrm{bd}}$ of degree $b$ (not in
general a commuting product of scalar factors), and $\Theta = \eta\mathbf SB_{\mathrm{bd}}$ is rational inner. Then
$\dim K = \deg\det\Theta = N$ (the resonances with multiplicity, delay modes included), $\mathrm{spec}\,Z$ is the
disc zero set of $\Theta$ with its Jordan structure, $Z^m\to0$, and
\[
I-Z^{*}Z = J^{*}J,\qquad \mathrm{rank}\,J\le h,\qquad \sum_{m\ge0}Z^{*m}J^{*}JZ^m = I ,
\]
so no nonzero vector has zero flux at every time (an individual instantaneous flux may vanish, ledger 23).
For eigenvectors $Zk_n = z_nk_n$ and exit amplitudes $a_n = Jk_n\in\mathbb C^h$,
$(1-\bar z_nz_m)\langle k_n,k_m\rangle = a_n^{*}a_m$, in particular $a_n\ne0$; in the scalar case
$k_a(w) = \Theta(w)/(w-a)$ has $Jk_a = 1$ and $\|k_a\|^2 = (1-|a|^2)^{-1}$. For $\mathrm{D2}$, with $\eta = -1$,
\[
\Theta_2(z) = \frac{z^2(2z^4-2z^2+1)}{z^4-2z^2+2},\qquad B_{\mathrm{bd}} = b_{1/\sqrt2}\,b_{-1/\sqrt2},\qquad N = 6 .
\]
Under H-LP this $K$, $Z$, $J$ is the geometric wave compression of \cref{thm:cusp-wave-contraction} with $h$ exits.
\end{theorem}

\begin{theorem}[The renewal channel with $h$ exits]
\label{thm:h-exit-renewal}
\claimstatus{thm:h-exit-renewal}{proved}
Under H-CONTRACTION, for every rebound map $\mathsf R(X) = \sum_k V_k XV_k^{*}$:
(a) $\mathcal E_{\mathsf R}$ is CPTP with Kraus operators $Z$ and $V_k J$.
(b) The holding instrument is an $h^2\times h^2$ transfer operator on exit matrices, not an $h\times h$
transition matrix (ledger 24); telescoping the defect identity gives total probability one, so $\mathcal R$ is
CPTP, and the mean holding time $\bar t(\mathsf u) = \sum_{k\ge0}\mathrm{Tr}\,\mathcal E_0^k\mathsf R(\mathsf u)$ is finite uniformly
in the exit density $\mathsf u$.
(c) With $\mathcal U = (I-\mathcal E_0)^{-1}\mathsf R$, the map $\mathsf u\mapsto\rdens_\infty = \mathcal U\mathsf u/\mathrm{Tr}\,\mathcal U\mathsf u
= \bar t(\mathsf u)^{-1}\sum_kZ^k\mathsf R(\mathsf u)Z^{*k}$ is a bijection between stationary exit densities of $\mathcal R$ and
stationary densities of $\mathcal E_{\mathsf R}$ (every stationary density has nonzero flux, by stability). Hence
uniqueness is uniqueness for $\mathcal R$, \emph{not} irreducibility (ledger 25: irreducibility is sufficient and gives a
faithful state; a unique absorbing state is a reducible example, \cref{prop:d2-canonical-rebound}); with
uniqueness, attraction of every density is absence of other peripheral eigenvalues, decided by the secular
identity $\det(I-u\mathcal E) = \det(I-u\mathcal E_0)\det(I-\widehat{\mathcal M}(u))$, $\widehat{\mathcal M}(u) =
\sum_mu^m\mathcal M_m$: a simple zero at $u = 1$ and no other circle zero of the second factor. A scalar gcd of
holding times does not replace this for a general instrument.
(d) For a \emph{fixed reset} the scalar proof of \cref{thm:cusp-renewal-discrete} applies for any exit rank:
the stationary density is unique, and for a modal $\reb = \sum_ip_i|e_i\rangle\langle e_i|$,
$\bar t = \sum_ip_i/(1-|z_i|^2)$, $\rdens_\infty = \bar t^{-1}\sum_i\frac{p_i}{1-|z_i|^2}|e_i\rangle\langle e_i|$, so every
modal reset on a family of eigenvectors is stationary iff their moduli coincide (\cref{prop:cusp-modal-rh}),
with holding law $(1-r^2)r^{2(m-1)}$ and mean $(1-r^2)^{-1}$; a modal reset charging both a delay mode and the
Hasse--Weil modes is never stationary. For an exit-dependent $\mathsf R$ this is false in both directions:
$Z = \mathrm{diag}(r_1,r_2)$, $J = \mathrm{diag}(\sqrt{1-r_1^2},\sqrt{1-r_2^2})$ and $\mathsf R = \mathrm{id}$ make every
diagonal density stationary with $r_1\ne r_2$ (ledger 27).
(e) A density is stationary for some rebound iff $\rdens-Z\rdens Z^{*}\ge0$ (then the fixed reset
$\reb = Q/\mathrm{Tr}\,Q$ works), and invariant flags realise every decreasing spectrum
(\cref{thm:cusp-admissible-flags}, arbitrary defect). On $\mathbb C\,\mathrm{vac}\oplus K$ with $\tilde Z = 1\oplus Z$,
$\tilde J = 0\oplus J$, the odd coherences $|k_n\rangle\langle\mathrm{vac}|$ are eigen-operators with eigenvalue $z_n$
for every $\mathsf R$ (\cref{prop:cusp-odd-protected}).
\end{theorem}

\begin{proposition}[Symmetry blocks; what the extra cusps add at genus one]
\label{prop:covariant-channel-blocks}
\claimstatus{prop:covariant-channel-blocks}{proved}
A finite automorphism group of the diagram acts unitarily on core, rays and exits; with a compatible inner
generator $Z$ commutes with the model representation and $J$ intertwines it with the exit representation, so
the \emph{no-event model} is an orthogonal sum of isotypic components. For $\mathrm{D3}$, $K = K_e\oplus K_o$,
$\dim K_e = 6$ (four Hasse--Weil modes and a length-two delay chain), $\dim K_o = 2$ (the delay chain of the
$z^2$ channel); the Dirichlet channel $-1$ has no model space. Within $K_e$ the Hasse--Weil and delay
subspaces are invariant but generally not orthogonal; the same holds for $\mathrm{D2}$. Covariance of a
\emph{channel} does not preserve these Hilbert-space blocks: it preserves isotypic components of the
conjugation representation on operator space, and densities of two inequivalent one-dimensional types both
transform trivially, so a covariant rebound may transfer populations between them (the brief's stronger claim
is false, ledger 29). At genus one the extra character channels add delay and constant factors and rebound
choices, no further nonzero resonances; they are not dynamically irrelevant.
\end{proposition}

\subsection{Which rebound does the arithmetic pick? The one-cusp cavity}

\begin{proposition}[The canonical cusp rebound of $\mathrm{D2}$ is an absorbing delay mode]
\label{prop:d2-canonical-rebound}
\claimstatus{prop:d2-canonical-rebound}{proved}
Since $\Theta_2(0) = 0$ (a consequence of $c = 1$), $j = J^{*}1 = \Theta_2(w)/w$ has $\|j\| = 1$ and $Zj = P_K\Theta_2 = 0$;
$g = \Theta_2(w)/w^2\in K$ is orthogonal to $j$, normalised, with $Zg = j$: the delay sector is one size-two Jordan
block. In an orthonormal basis (polynomial basis $w^k/(w^4-2w^2+2)$, Cholesky),
\[
Z = \begin{pmatrix}
0&2/3&0&-\sqrt5/6&0&0\\ 1&0&0&0&0&0\\ 0&\sqrt5/3&0&1/3&0&0\\ 0&0&1&0&0&0\\ 0&0&0&\sqrt3/2&0&0\\ 0&0&0&0&1&0
\end{pmatrix},\qquad j = e_6,\qquad I-Z^{*}Z = e_6e_6^{*},\qquad \chi_Z = z^2(z^4-z^2+\tfrac12).
\]
The canonical rebound $\reb_J = J^{*}J/\mathrm{Tr}(J^{*}J) = |j\rangle\langle j|$ (``flux returns where it left'') has
$\rdens_\infty = \reb_J$, spectrum $(1,0,0,0,0,0)$, holding law $m(m) = \delta_{m1}$, $\hat m(u) = u$, mean $1$: it is
the unique stationary density (fixed reset), attracting, and a \emph{modal rebound at zero} whose Hasse--Weil
Riesz projection vanishes although its overlaps $\langle j,k_a\rangle = 1$ with every Hasse--Weil mode are nonzero. It
never charges the zeros; the naive ``arithmetic picks the rebound'' candidate fails \emph{for this cut}: the degeneracy
is a property of where the core/cusp cut is placed (\cref{prop:multi-exit-spectrum}(iii), review finding 2), not of the
diagram. (On APW's convention with
the delay modes dropped it is a genuine mixture: law $(3/4,0,1/12,0,1/48,\dots)$ on odd times, mean $7/3$,
spectrum $\{1/7,1/7,1/7,4/7\}$; numerics D-ledger.) The same-defect amplitude is $\langle j,(I-uZ)^{-1}j\rangle = 1$,
whereas the Sz.-Nagy--Foias transfer $\langle k,[-Z+uD_{Z^{*}}(I-uZ^{*})^{-1}D_Z]j\rangle = \Theta_2(u)$ links
\emph{opposite} defect vectors, $I-ZZ^{*} = kk^{*}$, $k = (\sqrt{15}/6,0,-1/\sqrt3,0,1/2,0)^T$: the
characteristic function does not generate the holding law (ledger 32). \emph{Correction to
\cref{thm:cusp-renewal-discrete}(c):} ``aperiodic iff $c\ne1$ when the range of $\reb$ contains an eigenvector'' is
wrong; every normalised eigenvector has one-step exit probability $1-|z_i|^2>0$, including $z_i = 0$, so every
modal reset is aperiodic even at $c = 1$ (ledger 43; the periodic examples are non-modal).
\end{proposition}

\begin{proposition}[Modal Hasse--Weil rebounds of $\mathrm{D2}$; the Gram matrix in closed form]
\label{prop:d2-modal-gram}
\claimstatus{prop:d2-modal-gram}{proved-conditional}
For any weights on the four normalised Hasse--Weil eigenvectors, $\rdens_\infty = \reb$, $m(m) = \delta\,2^{-(m-1)/2}$,
$\bar t = \delta^{-1} = 2+\sqrt2$, $\delta = 1-1/\sqrt2$ (in general $(1-q^{-1/2})^{-1}$; the modular surface's $2$ is
continuous time), and the exit rate $\langle j|\reb|j\rangle = 1-q^{-1/2}$ for every modal reset, forced by the Gram
identity. With $a = \sqrt{(1+i)/2}$ (positive real and imaginary parts), roots ordered $(a,-a,\bar a,-\bar a)$,
$k_i = \Theta_2/(w-z_i)$, $e_i = \sqrt\delta\,k_i$, the normalised Gram matrix is
\[
O = \begin{pmatrix}1&d&u&v\\ d&1&v&u\\ \bar u&\bar v&1&d\\ \bar v&\bar u&d&1\end{pmatrix},\qquad
d = 3-2\sqrt2,\quad u = \delta(1-i),\quad v = \delta(3+i)/5 ,
\]
invariant under the parity matrix $\mathsf P$ swapping $1\leftrightarrow2$, $3\leftrightarrow4$, with
$\mathsf P\,\mathrm{diag}(z_i)\mathsf P = -\mathrm{diag}(z_i)$. The stationary density's spectrum is
$\mathrm{spec}(\mathrm{diag}(\sqrt p)\,O\,\mathrm{diag}(\sqrt p))$, the Gram spectrum, not the weight list: for equal
weights $0.11448581, 0.16190739, 0.29972775, 0.42387905$ (plus two zeros on the full $K$).
\end{proposition}

\begin{theorem}[The metrics in which $q^{1/4}Z$ is unitary; the arithmetic metric needs one exit per mode]
\label{thm:arithmetic-metric}
\claimstatus{thm:arithmetic-metric}{proved-conditional}
On $K_{\mathrm{HW}}$ of $\mathrm{D2}$, in the eigenbasis, a positive matrix $H$ defines an inner product for which
$Z^{*}HZ = r^2H$, $r = 2^{-1/4}$, iff $(\bar z_iz_j-r^2)H_{ij} = 0$ iff $H = \mathrm{diag}(h_1,\dots,h_4)$, $h_i>0$: a
cone of real dimension four, three modulo scale. Parity invariance forces $h_1 = h_2$, $h_3 = h_4$; invariance under
the real structure $z\mapsto\bar z$ forces $h_1 = h_3$, $h_2 = h_4$; both leave a \emph{single ray}, all weights
equal on the vectors $e_i$ of \cref{prop:d2-modal-gram} (the two symmetries act transitively on the four modes).
The Lax--Phillips energy metric is outside the cone: $|\langle e_a,e_{-a}\rangle| = 3-2\sqrt2$,
$|\langle e_a,e_{\bar a}\rangle| = \sqrt2-1$, $|\langle e_a,e_{-\bar a}\rangle| = (\sqrt2-1)/\sqrt5$, angles
$80.12^\circ$, $65.53^\circ$, $79.32^\circ$, so $q^{1/4}Z$ is not energy-unitary. This is the notebook's ``RH =
Ramanujan = the compressed operator is $q^{-1/4}$ times a unitary'' made literal for the elliptic cavity: true for
the cone, canonically for its symmetric ray, false for the energy metric; no independent global Ramanujan
statement follows. The channel of \cref{thm:h-exit-renewal} uses the energy metric and its adjoints; changing
metric rebuilds adjoints, defects and resets (ledger 35). Obstruction: in an arithmetic metric
$I-Z^{*'}Z = (1-r^2)I$ on $K_{\mathrm{HW}}$ has rank \emph{four}; it cannot factor through the one-dimensional cusp
exit, whose defect restricted to $K_{\mathrm{HW}}$ has rank one. An arithmetic-metric renewal construction needs at
least four exit coordinates on $K_{\mathrm{HW}}$, one per zero mode: a different conservative completion of the same
eigenvalue data (ledger 44).
\end{theorem}

\begin{observation}[H-HECKE: Hecke and Bost--Connes selection of the rebound is open]
\label{obs:hecke-rebound-open}
\claimstatus{obs:hecke-rebound-open}{open}
The stabiliser sizes specify adjacency, not the vertex and edge groups, their embeddings, or the double cosets
that define Hecke correspondences, so no Hecke operator can be built from H-DIAG alone. Commutation with
adjacency and preservation of cusp asymptotics would not by themselves force a modal rebound (joint eigenspaces
may coincide, $Z^2$ already has doubled eigenvalues, covariance of a CP map does not force diagonal outputs). A
selection theorem needs a specified Hecke action compatible with adjoints, simple separated joint characters and
a dephasing condition; the function-field Bost--Connes systems (Jacob, Consani--Marcolli, Neshveyev--Rustad;
\texttt{sources.md}) supply none of this for the diagram. \textbf{H-HECKE.}
\end{observation}

\subsection{The vacuum: finite volume, and funnels}

\begin{proposition}[The constant function on a finite-volume regular diagram never leaks]
\label{prop:constant-mode-no-exit}
\claimstatus{prop:constant-mode-no-exit}{proved}
On a graph with stabiliser function of finite volume ($\sum_v1/S(v)<\infty$) and bounded degree the constant
function is in $L^2(V,1/S)$ and $A$ is bounded self-adjoint; $A\mathbf 1 = \deg$, so \emph{regularity is
necessary} (ledger 37), and for a $(q+1)$-regular diagram $\mathbf 1$ is an eigenfunction with normalised
eigenvalue $(q+1)/\sqrt q>2$ and cusp tails $\propto q^{-k/2}$: a visible bound state with poles $z = \pm q^{-1/2}$
(bipartite case), not a cusp form. Under the discrete wave group $W(u,v) = (Tu-v,u)$ its data have time
multipliers $\sqrt q$ and $q^{-1/2}$, neither of modulus one, and the bound wave block has indefinite energy: the
Perron mode is \emph{not} stationary under the wave group and radiates no flux (\cref{obs:cusp-no-graph-vacuum}
holds, not an exception to it; a Schrödinger eigenstate or the Markov operator $A/(q+1)$ are other dynamics).
In the model with the bound data removed and the vacuum adjoined, $\tilde Z = 1\oplus Z$, $\tilde J = 0\oplus J$, every
rebound leaves the vacuum without an exit (\cref{prop:graded-stationary-segment}), independently of the number of
cusps: the hedgehog does not help. Finite volume does not forbid a different dynamics or a different defect on the
bound data.
\end{proposition}

\begin{proposition}[The funnel self-energy; the constant mode on a funnel]
\label{prop:funnel-self-energy}
\claimstatus{prop:funnel-self-energy}{proved}
APW's outgoing constant-type conditions are $f(w) = \sqrt q\,f(v)/\mathfrak m$ on a cusp and $f(w) = f(v)/(\sqrt q\,\mathfrak m)$ on a
funnel (\cref{cit:apw-resonance-matrix}, \texttt{:1915}), so after the notebook's normalisation the self-energies
are $\varSigma_{\mathrm{cusp}} = C_{\mathrm{cusp}}/z$ and $\varSigma_{\mathrm{funnel}} = C_{\mathrm{funnel}}/(qz)$ and
\[
\boxed{\ p_{\mathrm{funnel}}(z) = \det\big((1+z^2)I-zT_X-C_{\mathrm{cusp}}-q^{-1}C_{\mathrm{funnel}}\big)\ },
\]
with the same Laurent factor to APW's $H$; the one-vertex $(q+1)$-funnel tree gives $p = z^2-q^{-1}$, APW's
$\pm q^{-1/2}$, and the numerics lane reproduces their modular-curve example ($\pm q^{1/2}$) for $q = 2,3,5$. This is
the constant-type finite resonance matrix, not a proof that a funnel's wave space is a single standard ray. On a
funnel the constant function is not $L^2$ (sphere growth $q^k$), so the bound-state argument of
\cref{prop:constant-mode-no-exit} lapses and the constant function becomes outgoing at $z = q^{-1/2}$, the same
number with its role flipped from pole to zero. But a globally constant function is not outgoing on a cusp and a
funnel at once: the cusp condition needs $z = \sqrt q$, the funnel condition $z = 1/\sqrt q$. ``Zeros from the cusp,
vacuum leak from the funnel'' is a research target (ledger 40), and by \cref{cit:apw-cusps-pic-no-funnels} a
cusp-plus-funnel diagram is \emph{not} the quotient of the tree by a lattice in $\mathrm{PGL}(2,F)$: it must be built by
hand.
\end{proposition}

\begin{observation}[Genus, class number, and the number-field hedgehog: the testable prediction]
\label{obs:genus-class-prediction}
\claimstatus{obs:genus-class-prediction}{open}
Under H-CLASS a curve of genus $g$ with one removed rational point and $h$ cusps has one character channel
carrying the zeta numerator of degree $2g$, hence $4g$ nonzero square-root resonances (bipartite pairing is the
square root, not a second doubling), and each nontrivial unramified character an $L$-polynomial of degree $2g-2$,
hence $4g-4$ resonances: at genus one none (as $\mathrm{D3}$ shows), at genus two eight zeta resonances and four per
nontrivial character, all on $|z| = q^{-1/4}$, $4h+4$ in total before delay, bound, cusp-form and threshold
bookkeeping. These are character channels, not individual geometric spikes. For number fields every class
character carries Hecke $L$-zeros (\cref{cit:sorensen-scattering-matrix}). The $\Gamma_0(N)$ direction, one cusp
per divisor of the level with Dirichlet $L$-functions, remains H-ARITH (\cref{obs:cusp-arith-next}); the level tower
is the honest ``cusp per prime'', and its diamond operators carry $(\mathbb Z/N)^\times$, whose inverse limit is the
Bost--Connes symmetry group.
\end{observation}

\begin{numerical}[Blind numerics lane, elliptic cavities]
\label{num:elliptic-cavity}
\claimstatus{num:elliptic-cavity}{numerical}
\texttt{scripts/elliptic\_cavity.py} ($192$ checks, $51$ s, seed $20260920$; \texttt{outputs/elliptic\_cavity.txt};
ledger D1--D35 in \texttt{notes/elliptic-cavity/numerics.md}): both figures re-read and confirmed vertex for vertex;
degrees $q+1$; $p,\tilde p$ by exact rational arithmetic in $t = z/\sqrt q$ (integer coefficients) against APW's
polynomials; $\mathbf S$ from a direct generalised-eigenfunction solve against the closed form (deviation
$4\cdot10^{-15}$), symmetry, unitarity, functional equation, $\det\mathbf S = (-1)^hp/\tilde p$, $\det H = (-1)^n(2\mathfrak m)^{-n}p$;
$1/R_2$ against $z^{-2}\zeta_K(2s-1)/\zeta_K(2s)$; the three $\mathrm{D3}$ channels and $\det\mathbf S_e$ against the
zeta ratio; the height renormalisation $a-b = -1$ unique among half-integers in $[-4,4]$; $\mathrm{ord}_0p = 2\dim\ker(I-C)$;
$\dim\ker T_X = 3$ against two cusp forms; regular thresholds, $\mathrm{tr}\,\mathbf S(\pm1) = 0$, Levinson windings $4,6$;
the model contraction with its Jordan block, $\reb_J$ absorbing, modal means $2+\sqrt2$, exit rate $1-q^{-1/2}$,
Gram entries $\{2+\sqrt2,2-\sqrt2,\sqrt2,2/\sqrt{10}\}$; the metric cone and its single symmetric ray; the product
identity $[z^m]p = -1$ on both diagrams; the funnel formula against APW's tree and modular-curve examples for
$q = 2,3,5$.
\end{numerical}

\begin{observation}[What the next round should do]
\label{obs:elliptic-cavity-next}
\claimstatus{obs:elliptic-cavity-next}{sketched}
(i) The cusp-plus-funnel toy by hand: one core, one cusp carrying a zeta factor, one funnel, and the question
whether the renewal channel then has a unique mixed stationary state with the pole even and the zeros odd, given
\cref{prop:funnel-self-energy}'s incompatibility of the two outgoing conditions for a constant. (ii) A genus-two
diagram (or the $\mathbb Z/2\times\mathbb Z/2$ curve of Lorscheid) to test \cref{prop:nonzero-root-product}'s coefficient
$-1$ against Weil and \cref{obs:genus-class-prediction}'s $4h+4$. (iii) $\Gamma_0(N)$ over $\mathbb F_q[T]$ with $N$ a
product of two primes: four cusps, one per divisor, Dirichlet $L$-functions (H-ARITH). (iv) The arithmetic-metric
completion with one exit per zero mode (\cref{thm:arithmetic-metric}): what geometric object has four cusps seeing
the four Hasse--Weil modes separately? (v) Promote \cref{prop:d3-group-matrix} to a statement for general
$\mathrm{GL}_2(R)\backslash\mathcal T$ and prove it. (vi) Construct the Lax--Phillips completion for $\mathrm{D2}$, or stop
quoting $\dim K$ geometrically: H-LP has been inherited unexamined through three shards (reviewer).
\end{observation}
